\chapter{绪论}
\label{chap:intro}

\section{研究背景}
\label{sec:bg}

线上智能体大语言模型的能力来自持续的强化学习后训练：每隔若干周，用新采集的交互数据对策略做一次更新。这构成一个\textbf{任务无关的持续强化学习}问题——没有清晰的任务边界，新旧能力在同一参数空间里竞争。与离线训练一次成型不同，持续学习场景下，模型必须在新数据上习得新行为的同时，尽量不遗忘此前已经掌握的能力。

在智能体（agent）、多轮交互的设定下，这一问题比传统持续学习更困难，原因在于两个相互独立的失败模式：

\begin{enumerate}
  \item \textbf{灾难性遗忘（catastrophic forgetting）}：在新任务上的强化学习更新会系统性侵蚀旧领域能力。这是神经网络参数共享带来的固有难题，在持续学习文献中被广泛研究。
  \item \textbf{Echo Trap}\upcite{B4}：多轮智能体强化学习中，策略多样性会发生坍缩——同一查询的组内 rollout 轨迹趋于一致，导致组内归一化所得优势趋近于零，训练停滞乃至发散。这是多轮智能体场景特有的失效模式，单轮强化学习研究中并不暴露。
\end{enumerate}

已有研究表明，强化微调（RFT）比监督微调（SFT）更不易遗忘\upcite{B2,C1}，但其缓解作用有限，遗忘仍然发生；而多轮智能体场景进一步带来了单轮研究无法暴露的失败模式。因此，如何在多轮智能体大语言模型上做持续强化学习而不发生灾难性遗忘，是一个亟需回答的问题。

\section{研究问题与动机}
\label{sec:motivation}

针对灾难性遗忘的主流解法是\textbf{经验回放}（experience replay）\upcite{A2}及其各种改进——优先级缓冲区\upcite{A9}、按任务分区缓冲区\upcite{A3}、基于核心集的压缩缓冲区\upcite{A8}等。这些方法共享一个隐含假设：被回放的与在线的经验携带着\textbf{忠实的学习信号}，即用于更新策略的梯度方向是正确的，遗忘只源于"新数据覆盖旧能力"这一参数层面的竞争。

本文主张，在持续\textbf{多轮智能体}强化学习中，这个假设会悄悄破裂，而且破裂发生在任何回放或正则化\textbf{之前}。智能体场景的主力算法 GRPO（Group Relative Policy Optimization）通过对同一查询的 $M$ 条 rollout 在组内归一化来估计每条样本的优势 $A_i$。这个估计\textbf{仅当} $M$ 条轨迹的差异纯粹来自策略采样的随机性时才是无偏的。任何其他方差来源都会把环境噪声注入 $A_i$，从源头污染学习信号。在多轮智能体场景下，有两个典型效应在源头破坏这一前提：

\begin{enumerate}
  \item \textbf{组内环境发散}。跨多查询会话，$M$ 个执行环境逐渐分叉；某个槽位可能并非因策略更差、而是因前一条查询留下的不利状态被惩罚，于是优势吸收了随轮次累积的环境噪声。
  \item \textbf{多轮前提漂移}。后续查询通常引用上一轮的执行结果（例如"第 3 页的数据错了"）。当后续查询在数据采集时被写死、而实际状态是策略相关的随机函数时，其前提仅以一个\textit{随策略变强而漂移}的概率成立——于是奖励了对不存在问题的幻觉式修复，或惩罚了如实回答"没有这个问题"。
\end{enumerate}

被污染的优势在下游无法修复：任何回放加权或 KL 锚定都救不回一个符号本就错了的梯度。这一观察把持续智能体强化学习重新表述为一个\textbf{两阶段}问题——先产出干净信号，再巩固它——并把本文的整个系统围绕这一主旨来组织。

\section{主要贡献}
\label{sec:contrib}

在"先产出干净信号、再巩固它"的单一主旨下，本文的贡献如下：

\begin{enumerate}
  \item \textbf{对持续智能体强化学习的源头级重新表述}。本文识别并形式化了 GRPO 学习信号的两类污染——组内环境发散与多轮前提漂移——它们发生在巩固\textit{之前}，且无法用回放或正则化修复（第~\ref{chap:setup} 章）。
  \item \textbf{一套产出干净信号的数据与 rollout 系统}。winner 同步会话调度在多查询会话中恢复组内位级一致起点，以消除组内环境发散；观察、出题、奖励三个协作智能体在线构造前提成立的多轮数据，根除前提漂移，并兼作一个跟随策略能力前沿的自适应课程（第~\ref{chap:method} 章 \S\ref{sec:rollout}--\ref{sec:usersim}）。
  \item \textbf{一套与干净信号匹配的巩固方法}。以零源码改动注入 GRPO 的四项持续学习目标、按能力分区的抗遗忘回放缓冲区、以及 U 形块级回放重加权方案（第~\ref{chap:method} 章 \S\ref{sec:obj}--\ref{sec:reweight}）。
  \item \textbf{系统实现与实验设计}。基于 verl 框架零侵入落地的端到端训练系统（第~\ref{chap:system} 章），以及在 27B 智能体、九个能力域、Pass$^3$ 评测基准上的逐组件受控消融实验设计（第~\ref{chap:exp} 章）。
\end{enumerate}

\textbf{关于叙事纪律的说明。} 上述贡献 2 和 3 刻意不作为一串独立技巧来呈现；每一项都是服务于"先产出、再巩固一个未被污染信号"这一单一主旨的手段。在消融实验（第~\ref{chap:exp} 章）中，每个子组件都必须对着这一主旨挣到自己的位置；任何在干净信号机制下不改善遗忘的子组件，将作为负面结果如实报告，而非强行辩护。

\section{论文组织结构}
\label{sec:structure}

本文余下部分组织如下。第~\ref{chap:related} 章综述持续学习、经验回放、智能体强化学习等相关工作并定位本文的位置。第~\ref{chap:setup} 章给出问题定义，形式化 GRPO 优势估计的两类信号污染并重新表述为两阶段问题。第~\ref{chap:method} 章是核心方法章，分别介绍持续学习目标、九桶回放缓冲区、U 形块重加权、沙箱 winner 同步会话、三智能体在线多轮构造。第~\ref{chap:system} 章给出基于 verl 的系统实现。第~\ref{chap:exp} 章给出评测基准、消融实验设计与分析维度。第~\ref{chap:conclusion} 章总结全文并展望未来工作。

\textbf{关于实验状态的说明。} 截至本文写作时，训练系统代码已全部完成并通过约 290 项单元测试，但受限于 64 卡 GPU 集群与真实模型权重的可用性，21 个消融实验尚未实际运行。因此第~\ref{chap:exp} 章呈现的是\textbf{实验设计与评测协议}，所有依赖训练结果的数字以 \todo{训练后回填} 标记，本文不编造任何实验结果。
